\chapter{Experimental Methods}
\label{ch:experimental}

\epigraph{%
  The quality of a biomechanical analysis is limited by the quality of the measurements,
  not the sophistication of the model.
}{}

\section{Motion Capture Systems}

\subsection{Marker-Based Systems}

Optoelectronic motion capture systems (Vicon, Qualisys, OptiTrack) track retroreflective
markers attached to the body surface. Standard configurations use 30--80 markers placed
on bony landmarks following established protocols (Plug-in-Gait, Cleveland Clinic, ISB).

Key specifications for biomechanical applications
\cite{chiari2005human,chiari2005human}:
\begin{itemize}
  \item \textbf{Sampling rate}: 100--500~Hz for human motion (higher for impacts)
  \item \textbf{Spatial accuracy}: $<1$~mm in calibrated volume
  \item \textbf{Marker size}: 9--14~mm diameter
  \item \textbf{Latency}: important for real-time applications ($<10$~ms achievable)
\end{itemize}

It is important to note that the instrument accuracy above refers to the camera
system alone. The dominant error in marker-based biomechanics is not camera noise
but \emph{soft tissue artifact}---the relative motion between skin markers and the
underlying bone---discussed in Section~\ref{sec:sta} below
\cite{chiari2005human,chiari2005human}.

\subsection{Markerless Systems}

Computer vision approaches (OpenPose, MediaPipe, DeepLabCut) estimate joint positions
from video without physical markers. Accuracy is currently 10--30~mm, insufficient for
clinical-grade inverse dynamics but useful for field measurements and large-scale studies.

\begin{lstlisting}[language=Python, caption={Motion capture data processing pipeline.}]
import numpy as np
from scipy.signal import butter, filtfilt

def process_mocap_data(
    marker_data: np.ndarray,
    sample_rate: float,
    cutoff_freq: float = 6.0,
    filter_order: int = 4
) -> tuple[np.ndarray, np.ndarray, np.ndarray]:
    """Process raw motion capture marker data.

    Parameters
    ----------
    marker_data : np.ndarray, shape (n_frames, n_markers, 3)
        Raw marker positions in meters.
    sample_rate : float
        Sampling frequency in Hz.
    cutoff_freq : float
        Low-pass filter cutoff frequency in Hz.
    filter_order : int
        Butterworth filter order.

    Returns
    -------
    tuple of np.ndarray
        (positions, velocities, accelerations)
        Each shape (n_frames, n_markers, 3).
    """
    # Design low-pass Butterworth filter
    nyquist = sample_rate / 2.0
    b, a = butter(filter_order, cutoff_freq / nyquist, btype='low')

    n_frames, n_markers, _ = marker_data.shape
    positions = np.zeros_like(marker_data)
    velocities = np.zeros_like(marker_data)
    accelerations = np.zeros_like(marker_data)

    dt = 1.0 / sample_rate

    for m in range(n_markers):
        for axis in range(3):
            # Filter position data
            positions[:, m, axis] = filtfilt(
                b, a, marker_data[:, m, axis]
            )

    # Central difference for velocities
    velocities[1:-1] = (positions[2:] - positions[:-2]) / (2 * dt)
    velocities[0] = velocities[1]
    velocities[-1] = velocities[-2]

    # Central difference for accelerations
    accelerations[1:-1] = (
        positions[2:] - 2 * positions[1:-1] + positions[:-2]
    ) / dt**2
    accelerations[0] = accelerations[1]
    accelerations[-1] = accelerations[-2]

    return positions, velocities, accelerations
\end{lstlisting}

\subsection{Soft Tissue Artifact}
\label{sec:sta}

Soft tissue artifact (STA) is the relative motion between skin-mounted markers
and the underlying bone caused by deformation of skin, subcutaneous fat, and muscle
during movement. STA is widely recognized as the dominant source of error in
marker-based motion capture, frequently exceeding the nominal camera resolution
by more than an order of magnitude \cite{chiari2005human,peters2010quantification}.

The magnitude of STA varies strongly with marker location, anatomical region,
task, and subject anthropometry. Representative values for the lower limb are
10--30~mm of skin-marker displacement relative to bone during walking and running,
with the thigh segment the worst offender and the shank the best
\cite{chiari2005human,peters2010quantification}. STA is not random noise: it is
task-correlated and heavily dependent on muscle contraction state, which means
low-pass filtering cannot remove it.

Practical mitigation strategies include:
\begin{itemize}
  \item \textbf{Rigid marker clusters} mounted on lightweight shells, which average
        out local skin motion over a neighborhood of markers.
  \item \textbf{Optimal pose estimation}, in which the rigid body pose of a segment
        is obtained by least-squares fitting of a cluster of markers rather than by
        differencing pairs of landmarks \cite{chiari2005human}.
  \item \textbf{Global optimization / inverse kinematics} with joint constraints,
        which enforces an articulated model and redistributes STA across the chain.
  \item \textbf{Biplanar fluoroscopy} or bone pins when the research question
        demands sub-millimeter bone-to-bone accuracy.
\end{itemize}

STA should always be reported as part of the uncertainty budget of any kinetic
analysis; inverse dynamics is sensitive to segment orientation errors, and
moderate STA can produce large errors in computed joint moments.

\subsection{Low-Pass Filter Cutoff Selection}
\label{sec:filter-cutoff}

Choosing the cutoff frequency for the low-pass Butterworth filter in
\texttt{process\_mocap\_data()} is not a matter of taste: a cutoff that is too low
attenuates the true signal (especially accelerations, whose energy extends to
higher frequencies than positions), while a cutoff that is too high leaves noise
that is amplified by numerical differentiation.

The standard objective procedure is \textbf{residual analysis}, due to Winter
\cite{Winter2009}. For a range of candidate cutoffs $f_c$, one computes
the root-mean-square residual
\[
  R(f_c) = \sqrt{\frac{1}{N}\sum_{i=1}^{N}\bigl(x_i - \hat{x}_i(f_c)\bigr)^2},
\]
where $x_i$ is the raw signal and $\hat{x}_i(f_c)$ is the filtered signal. Plotting
$R(f_c)$ vs.\ $f_c$ produces a curve that is approximately linear in the high-$f_c$
regime (residuals reflect only noise) and rises sharply at low $f_c$ (residuals
include signal). The optimal cutoff is the intersection of the extrapolated linear
noise-only segment with the $R$-axis, typically in the range 4--10~Hz for
kinematics of whole-body locomotion and 10--20~Hz for running and impact-dominated
tasks \cite{Winter2009}. Higher cutoffs are required when downstream
analysis needs accelerations or joint powers, since differentiation acts as a
high-pass amplifier.

\section{Electromyography (EMG)}

EMG measures the electrical activity of muscles, providing a window into neural
commands. Surface EMG (sEMG) is non-invasive; intramuscular EMG (using fine-wire
or needle electrodes) targets deep muscles.

\subsection{EMG Processing}

The raw EMG signal must be processed before use:
\begin{enumerate}
  \item \textbf{Band-pass filter}: 20--450~Hz (remove movement artifact and noise)
  \item \textbf{Full-wave rectification}: $|e(t)|$
  \item \textbf{Envelope extraction}: low-pass filter at 3--6~Hz
  \item \textbf{Normalization}: divide by maximum voluntary contraction (MVC) value
\end{enumerate}

\section{Force Plates}

Force plates measure ground reaction forces (GRF) and moments. Standard specifications:
\begin{itemize}
  \item 6-component measurement: $F_x, F_y, F_z, M_x, M_y, M_z$
  \item Sampling rate: 1000--2000~Hz
  \item Range: 0--5000~N vertical, 0--2500~N horizontal
  \item Center of pressure (CoP) accuracy: $<2$~mm
\end{itemize}

The center of pressure is computed from force plate outputs:
\begin{equation}
  x_{\text{CoP}} = -\frac{M_y + F_x \cdot d_z}{F_z}, \quad
  y_{\text{CoP}} = \frac{M_x - F_y \cdot d_z}{F_z},
  \label{eq:ch7:cop}
\end{equation}
where $d_z$ is the distance from the force plate surface to the transducer plane.

\section{Inertial Measurement Units (IMUs)}

IMUs combine accelerometers, gyroscopes, and magnetometers in a compact package.
They provide:
\begin{itemize}
  \item Angular velocity: $\bm{\omega}(t) \in \Reals^3$
  \item Linear acceleration: $\bm{a}(t) \in \Reals^3$
  \item Magnetic field: $\bm{B}(t) \in \Reals^3$ (for heading)
\end{itemize}

Orientation estimation from IMU data is a state estimation problem solvable using
the extended Kalman filter (Volume~I, Chapter~6) or complementary filters.

\section{Data Synchronization}

Multiple measurement systems must be temporally synchronized. Best practices:
\begin{itemize}
  \item Hardware sync: trigger all systems from a common clock pulse
  \item Minimum: force plate and motion capture on the same clock
  \item EMG-to-motion delay: typically 10--80~ms (electromechanical delay)
  \item Cross-correlation for post-hoc alignment when hardware sync is unavailable
\end{itemize}

\section*{Exercises}

\begin{enumerate}
  \item \textbf{Filter Selection.}
        Using the provided code, filter a synthetic marker trajectory with cutoff
        frequencies of 4, 6, 8, and 12~Hz. Plot the effect on position, velocity,
        and acceleration. Discuss the trade-off between noise reduction and signal distortion.

  \item \textbf{CoP Computation.}
        Given force plate data from a vertical jump ($F_z$ peak = 3000~N), compute
        the center of pressure trajectory during takeoff.

  \item \textbf{EMG Processing.}
        Process a synthetic EMG signal (sum of sinusoids + white noise) through the
        full processing pipeline. Plot each stage.

  \item \textbf{(Programming.)} Implement an IMU orientation estimator using a
        complementary filter. Test with synthetic gyroscope and accelerometer data.
\end{enumerate}
